\documentclass{article}%
\usepackage[T1]{fontenc}%
\usepackage[utf8]{inputenc}%
\usepackage{lmodern}%
\usepackage{textcomp}%
\usepackage{lastpage}%
\usepackage{geometry}%
\geometry{margin=2.2cm}%
%
\usepackage{amsmath}%
\usepackage{amssymb}%
\usepackage{booktabs}%
\usepackage{longtable}%
\usepackage{array}%
\usepackage{hyperref}%
\hypersetup{hidelinks}%
\title{Soluções\\[4pt]\large Comparação com o gabarito Gabarito oficial AFA 2027, Versão A\thanks{Gerado com gpt-resolve: \url{https://github.com/lgabs/gpt-resolve}}}%
\date{29/10/2026}%
%
\begin{document}%
\normalsize%
\maketitle%
\section*{Resumo}\textbf{Nota média: 9,4/10} (15 de 16 itens com a resolução correta; cada item vale o mesmo dentro da questão).\par\medskip%
\begin{longtable}{@{}c c p{0.78\textwidth}@{}}\toprule Questão & Nota & Resumo \\ \midrule
17 & 10,0 & A resolução e o gabarito batem e estão certos: a alternativa é b) e só a afirmativa II é verdadeira. Pela figura original, f atinge 3k uma única vez, no ramo que vai a +$\infty$ perto de 6a, porque o ponto (-2b, 3k) tem bolinha aberta. A I é falsa porque -a pertence ao domínio. A III é falsa porque f não tem máximo. A IV é falsa porque a imagem é IR. A V é falsa porque f decresce em {[}-4b, -3b{]}. A VI é falsa porque f não é injetora. Os argumentos da resolução estão certos e completos, e não achei erro de leitura do gráfico. \\
18 & 10,0 & A resolução está certa e bate com o gabarito oficial: alternativa D, A = {]}-4, 2{]} ? {[}3, 4{]} ? \{5/2\}. Conferi a transcrição com o PDF original e está fiel. O domínio de f, a análise do sinal do logaritmo nos dois casos, a inequação de g, a interseção e o complementar estão corretos. Fiz uma verificação numérica que confirma o conjunto inteiro, extremidades incluídas. Não há divergência nem ponto duvidoso. \\
19 & 10,0 & A resolução e o gabarito oficial da Versão A coincidem na alternativa (A). Refiz o domínio e ele é M = {]}2k$\pi$, $\pi$/3 + 2k$\pi${[} (sen x > 0 e cos x > 1/2). O intervalo {[}$\pi$/6, $\pi$/4{]} está contido em M, e as justificativas para descartar (b), (c) e (d) estão certas. Há só um deslize de redação ('três condições', quando são duas), sem impacto na resposta. \\
20 & 10,0 & A resolução e o gabarito concordam: a resposta da questão 20 é a alternativa D, e ela está correta. O prisma tem aresta da base 11 e diagonal 11$\surd$11, logo altura 33. Ele passa em pé pela abertura de 16$\times$16 e cabe exatamente na altura da caixa. As outras três alternativas falham. A esfera tem diâmetro 32, maior que 20. A esfera inscrita no cone tem diâmetro 18, maior que a abertura de 16. O octaedro tem aresta 40, e sua esfera inscrita tem diâmetro de cerca de 32,7, maior que 20. Todas as contas foram conferidas, e o argumento da esfera inscrita está correto. \\
21 & 10,0 & A resolução e o gabarito oficial concordam: alternativa B, área 2000. A identificação da hipérbole (a$^2$ = 420, b$^2$ = 80, c$^2$ = 500) está certa, a posição de D e E no eixo y está bem justificada, e a área do losango vale 4c$^2$ = 2000. Conferi tudo de forma independente e não encontrei erro. \\
22 & 0,0 & A resolução marca B e o gabarito oficial marca D; quem está certo depende de como se lê a afirmativa (II). As contas da resolução estão corretas e eu as conferi. As raízes são 2, -1, $\pm$5i. (I) é verdadeira em qualquer convenção de argumento principal, porque os argumentos formam uma PA de razão $\pi$/2. (III) também é verdadeira: $\sigma$$_{2}$ = 23, que é primo. Em (II), lida ao pé da letra ('há um polígono'), o triângulo de vértices -1, 5i e -5i tem área 5. Ela fica verdadeira e a resposta é B. Lida como 'o polígono dos quatro afixos', a área é 15, que é múltiplo e não divisor de 5. Ela fica falsa e a resposta é D, que parece ser o que a banca pretendia. O enunciado é ambíguo e a questão merece recurso, por isso marquei 'unclear'. \\
23 & 10,0 & A resolução está correta e bate com o gabarito oficial (B). As horas são Caravan 52 h, Phenon 32 h e Eurocopter 71,5 h. Lendo os custos pelo nome da aeronave, x = 4 + 3 + 4 = 11 mil dólares, logo 9 $\leq$ x < 13. Conferi as contas em Python e a transcrição com o PDF original. Não há erro na resolução nem no gabarito. \\
24 & 10,0 & A resolução e o gabarito oficial coincidem na alternativa A, 2 + P\_3$\cdot$C\_\{26,3\}$\cdot$A\_\{10,2\}$\cdot$C\_\{21,1\} = 29 484 002. Conferi o enunciado no PDF e as contas em python, e está tudo correto. A soma (e não o produto) dos três modos exclusivos de entrar elimina c) e d), que além disso têm o mesmo valor numérico entre si. O fator P\_3, que separa a) de b), vem da ordem livre dos blocos, e essa leitura é a confirmada pelo gabarito. As únicas ressalvas são ambiguidades de redação do enunciado (ordem dos blocos, \textbackslash{}"pares ordenados\textbackslash{}" no plural), que não mudam a nota. \\
25 & 10,0 & A resolução e o gabarito batem: os dois dão a alternativa A, 8$\pi$(2 + $\surd$3/3) $\approx$ 64,8 u.v. Refiz a fatoração de P e a geometria do trapézio. Calculei o volume de três formas (cilindro menos dois cones, Pappus--Guldin e integração numérica por arruelas) e todas confirmam o valor. A resolução está correta e bem justificada, sem erros ou ambiguidades relevantes. \\
26 & 10,0 & Questão 26: a resolução e o gabarito oficial (Versão A) dão a mesma resposta, (D), e ela está certa. A população fica entre 0 e 70 mil, cai com o tempo e tem inflexão em (10, 35) mil, ou seja, P(t) = 70/(1+e\^{}\{k(t-10)\}). A dedução das propriedades da logística (imagem, inflexão e simetria) está certa. O ponto que vale levar a um professor é o enunciado: ele diz \textbackslash{}"{]}-L, L{[}\textbackslash{}", o que contradiz a fórmula e o exemplo, e lido ao pé da letra levaria à alternativa (a). Isso pode motivar recurso, mas a resposta mais coerente continua sendo (D). \\
27 & 10,0 & A resolução está correta e coincide com o gabarito oficial (alternativa B). Os valores são Q1 = 1, Q2 = 3(3$\surd$3-$\pi$)/(9$\surd$3-$\pi$) $\approx$ 0,495 e Q3 = (6$\surd$3-3$\pi$)/(6$\surd$3-2$\pi$) $\approx$ 0,235, logo Q3 < Q2 < Q1. Conferi a leitura das áreas cinza na figura original ampliada e refiz as contas numericamente. Os raios $\rho$ = R/3 (círculos dos vértices do triângulo) e $\rho$' = a/3 (círculos do hexágono) estão bem justificados. Não há erro nem no gabarito nem na resolução. \\
28 & 10,0 & A resolução está correta e coincide com o gabarito oficial: alternativa (A). A condição de ao menos 3003 tiros dá 5x + 7y - 91 $\geq$ 0, e o orçamento de no máximo R\$ 26000 dá 33x + 242y - 1300 $\leq$ 0. Todas as contas e as análises das alternativas conferem. O único ponto a levar ao professor é uma ambiguidade do enunciado, que a resolução já aponta: lendo 'necessariamente' ao pé da letra, os pares que atendem às condições também satisfazem (b). A intenção da banca, porém, é claramente a (a). \\
29 & 10,0 & Questão 29: a resolução e o gabarito oficial dão a mesma resposta, alternativa C. Conferi a leitura da figura no PDF original. A lei f(x) = 3 + 2cos($\pi$/3$\cdot$(x - 9/4)) passa por todos os pontos marcados, e f(m) - m = f(3) - 3 = $\surd$2 $\approx$ 1,41 $\in$ {[}1, 3/2{[}. A resolução está correta e justifica cada passo, inclusive por que o resultado não depende do sinal de t e q. Não há erro nem no gabarito nem na resolução. \\
30 & 10,0 & A questão foi anulada pela banca, que aceita qualquer alternativa, e por isso a resposta C) 11 da resolução está de acordo com o gabarito. Refazendo as contas, h?(6) e j?(1) estão certos. Já g = -c$\cdot$f tem intercepto -c$^2$ e vértice com ordenada -c$\cdot$y\_v, enquanto o gráfico (4) marca -c e -y\_v, que são valores de -f. Com c > 1, nenhum gráfico representa g exatamente, o que justifica a anulação. A resolução notou essa diferença e escolheu a resposta que a banca provavelmente queria (4 + 6 + 1 = 11). A anulação e a escolha da resolução se sustentam. Não salvei nenhuma memória porque o acesso à pasta de memória foi bloqueado. \\
31 & 10,0 & A resolução e o gabarito dão a mesma resposta: alternativa c), aeronave 5, e ela está correta. A média é 6 h e o desvio padrão populacional é $\surd$6,5 $\approx$ 2,55 h. O valor mais próximo é 3 h, da aeronave 5 (distância 0,45, contra 0,55 até as 2 h da aeronave 8). A conclusão vale também para o desvio padrão amostral ($\approx$ 2,73 h). A resolução está correta, completa e bem justificada. \\
32 & 10,0 & A resolução e o gabarito concordam: resposta (C) 4. Pelo complementar, 1 - (2/5)\^{}n > 0,95 dá n > 3,25 com os logaritmos dados, e a conta direta confirma: n = 3 dá 93,6\% e n = 4 dá 97,44\%. A resolução está certa e bem justificada. A inversão do sinal ao dividir pelo log negativo está explícita, e a hipótese de independência (rebanho grande) foi declarada. Só há uma ressalva menor sobre essa hipótese de modelagem. \\
\bottomrule\end{longtable}%
\section*{Principais pontos de dúvida}%
\begin{itemize}%
\item%
\textbf{Q22} \textit{(statement)} A afirmativa (II) é ambígua. 'Considerando os afixos de x1, x2, x3 e x4, há um polígono...' pode querer dizer (a) o polígono cujos vértices são os quatro afixos (quadrilátero, área 15, não divide 5, logo falsa e gabarito D) ou (b) existe algum polígono com vértices entre os afixos (o triângulo -1, 5i, -5i tem área 5, logo verdadeira e resposta B). O 'há um' favorece (b). A pegadinha divisor $\times$ múltiplo (15 é múltiplo de 5) indica que a banca pensou em (a). Pode ser motivo de recurso, anulação ou troca de gabarito; vale ver se a banca publicou o gabarito definitivo depois dos recursos.%
\item%
\textbf{Q22} \textit{(key)} O gabarito D só se sustenta se (II) for falsa. (I) é verdadeira tanto em {[}0,2$\pi${[} quanto em {]}-$\pi$,$\pi${]}, e (III) dá $\sigma$$_{2}$ = 23, primo. Então o gabarito depende de ler (II) como 'o quadrilátero formado pelos quatro afixos', o que não está escrito no enunciado. Pelo jeito, a banca não percebeu que o triângulo -1, 5i, -5i tem área exatamente 5.%
\item%
\textbf{Q26} \textit{(statement)} O enunciado afirma que "{]}-L, L{[} é o intervalo de mínimo e máximo" da curva. Isso contradiz a própria lei, cuja imagem é {]}0, L{[}, e o exemplo, em que L = 6 e a curva fica entre 0 e 6. Lendo ao pé da letra, amplitude 70 mil daria 2L = 70, ou seja, L = 35, e uma curva decrescente entre 35 e -35, que é exatamente o gráfico (a). A alternativa (a) parece ter sido feita para quem segue essa frase. Essa incoerência interna pode motivar recurso ou anulação. Mesmo assim, a fórmula, o exemplo e o contexto (população não pode ser negativa) levam a (d), como dizem a resolução e o gabarito.%
\item%
\textbf{Q30} \textit{(statement)} Pelo enunciado, g(x) = -c$\cdot$f(x) com c > 1. Então g(0) = -c$^2$ e o vértice de g é (x\_v, -c$\cdot$y\_v). O gráfico (4), o único com concavidade para cima e raízes x$_{1}$ e x$_{2}$, marca -c e -y\_v, que valem para -f(x). Nenhum dos gráficos representa g de forma exata, e esse é muito provavelmente o motivo da anulação. A resolução afirma que '(4) é o único gráfico compatível'. Isso força a interpretação, porque com esses rótulos o gráfico (4) não é o de g.%
\end{itemize}%
\section*{Questão 17}%
\begin{longtable}{@{}p{0.06\textwidth} p{0.33\textwidth} p{0.33\textwidth} p{0.1\textwidth} p{0.08\textwidth}@{}}\toprule Item & Resolução & Gabarito & Veredito & Certo \\ \midrule
17 & (B) apenas uma afirmativa é verdadeira (só a II) & B & igual & ambos \\
\bottomrule\end{longtable}%
\begin{itemize}%
\item%
\textbf{17)} Refiz a análise sobre a figura original. (I) é falsa porque há bolinha fechada em (-a, 5k/2), então -a $\in$ A e A = {]}-5b, 6a{[}. (II) é verdadeira: o valor 3k não é atingido no ramo {]}-2b, -a{[}, já que (-2b, 3k) tem bolinha aberta e f(-2b) = 2k. Também não é atingido na poligonal (f $\leq$ 2k), no patamar (5k/2) nem em {[}a, 5a/2{]} (f $\leq$ 5k/2). Em {]}5a/2, 6a{[}, f cresce de -k até +$\infty$, então atinge 3k uma única vez. (III) é falsa porque f $\to$ +$\infty$ quando x $\to$ 6a$^{-}$. (IV) é falsa porque Im = {]}-$\infty$, 3k{[} ? {[}-k, +$\infty${[} = IR. (V) é falsa porque f decresce em {[}-4b, -3b{]}. (VI) é falsa porque f não é injetora (f(-4b) = f(-2b) = 2k, e o patamar em {[}-a, a{]}), então f$^{-}$$^{1}$ não existe. Só uma afirmativa é verdadeira, alternativa b). A resolução e o gabarito batem e estão certos.%
\end{itemize}%
\paragraph{Pontos de dúvida}%
\begin{itemize}%
\item%
\textit{(statement)} A resposta depende de ler na figura se as bolinhas estão abertas ou fechadas. Se (-2b, 3k) fosse fechada, 3k seria atingido duas vezes e a II ficaria falsa (resposta a). Se (-a, 5k/2) fosse aberta, a I ficaria verdadeira. No original, (-2b, 3k) está aberta e (-a, 5k/2) e (-2b, 2k) estão fechadas, o que confirma a leitura da resolução.%
\item%
\textit{(solution)} A resolução trata o trecho {[}5a/2, 6a{[} como estritamente crescente, inclusive perto de (4a, k/2), onde a curva quase se achata (parece haver um ponto de inflexão com tangente horizontal). Essa leitura do gráfico é razoável, e uma tangente horizontal isolada não tira a monotonia estrita. Mesmo assim, é uma hipótese tirada da figura.%
\item%
\textit{(statement)} A afirmativa VI é falsa porque f não tem inversa, mas a redação tenta o aluno a pensar só na reflexão em torno de y = x (x = -a $\to$ y = -a, x = 6a $\to$ y = 6a). Quem aceitasse essa reflexão contaria duas verdadeiras e marcaria d). A resolução trata esse ponto corretamente.%
\end{itemize}%
\section*{Questão 18}%
\begin{longtable}{@{}p{0.06\textwidth} p{0.33\textwidth} p{0.33\textwidth} p{0.1\textwidth} p{0.08\textwidth}@{}}\toprule Item & Resolução & Gabarito & Veredito & Certo \\ \midrule
18 & (D) A = {]}{-}4, 2{]} ? {[}3, 4{]} ? \{5/2\} & D & igual & ambos \\
\bottomrule\end{longtable}%
\begin{itemize}%
\item%
\textbf{18)} Refiz a questão. Para f estar definida, é preciso x $\neq$ 2, 5/2, 3, 4. A condição f > 0 vale quando base e logaritmando estão do mesmo lado de 1. Com os dois maiores que 1, sai x < 2 ou x > 4. Com os dois entre 0 e 1, sai {]}2,3{[} - \{5/2\}. A condição g $\leq$ 0 equivale a (x+4)(x-2) $\geq$ 0, ou seja, x $\leq$ -4 ou x $\geq$ 2. Então S = {]}-$\infty$,-4{]} ? {]}2,5/2{[} ? {]}5/2,3{[} ? {]}4,+$\infty${[}, e o complementar é A = {]}-4,2{]} ? \{5/2\} ? {[}3,4{]}, que é a alternativa (d). Conferi também em Python, numa malha de passo 0,001 em {[}-10, 10{]} que inclui os pontos -4, 2, 5/2, 3 e 4: x está fora de S exatamente quando está no conjunto da alternativa D, sem nenhuma exceção. Os valores numéricos citados na resolução também estão certos: g(-4) = 0, g(-10) = -216, g(-2) = 24, f(-4) = log$_{7}$13 e f(2,8) = log$_{0}$,$_{2}$0,6 > 0. Os contraexemplos usados para descartar a), b) e c) são válidos.%
\end{itemize}%
\section*{Questão 19}%
\begin{longtable}{@{}p{0.06\textwidth} p{0.33\textwidth} p{0.33\textwidth} p{0.1\textwidth} p{0.08\textwidth}@{}}\toprule Item & Resolução & Gabarito & Veredito & Certo \\ \midrule
19 & Alternativa (A): {[}$\pi$/6, $\pi$/4{]} ? M, com M = ? {]}2k$\pi$, $\pi$/3 + 2k$\pi${[} (k $\in$ $\mathbb{Z}$) & A & igual & ambos \\
\bottomrule\end{longtable}%
\begin{itemize}%
\item%
\textbf{19)} Conferi a transcrição com a página 8 do PDF e ela está fiel. Refiz o domínio. A segunda raiz exige sen x $\neq$ 0 e (1 - cos x)/sen x $\geq$ 0. Como 1 - cos x só se anula quando sen x = 0, isso dá sen x > 0. Com sen x > 0, a primeira raiz exige cos x > 1/2. Logo M = {]}0, $\pi$/3{[} + 2k$\pi$. O intervalo {[}$\pi$/6, $\pi$/4{]} está dentro de {]}0, $\pi$/3{[}. Em (b), o ponto $\pi$/2 dá radicando -2. O intervalo de (c) contém $\pi$/3, onde o denominador se anula. Em (d), cos x < 1/2 em todo o intervalo, então o radicando é negativo. Também conferi numericamente com Python, usando 20001 pontos por intervalo: só (a) fica inteiramente em M, e (d) não tem nenhum ponto em M.%
\end{itemize}%
\paragraph{Pontos de dúvida}%
\begin{itemize}%
\item%
\textit{(solution)} Problema de redação: a resolução diz que são 'três condições', mas lista só duas, (I) e (II). Isso não afeta o raciocínio nem a resposta.%
\item%
\textit{(statement)} O enunciado não define M. A resolução supõe, de forma razoável, que M é o domínio natural de g (o maior subconjunto de $\mathbb{R}$ onde g está definida). Essa é a leitura usual nas provas da AFA.%
\end{itemize}%
\section*{Questão 20}%
\begin{longtable}{@{}p{0.06\textwidth} p{0.33\textwidth} p{0.33\textwidth} p{0.1\textwidth} p{0.08\textwidth}@{}}\toprule Item & Resolução & Gabarito & Veredito & Certo \\ \midrule
20 & D & D & igual & ambos \\
\bottomrule\end{longtable}%
\begin{itemize}%
\item%
\textbf{20)} Conferi o enunciado no PDF e a transcrição está fiel. O interior da caixa mede 50$\times$20$\times$33 e a abertura mede 16$\times$16 (50-17-17 por 20-2-2). Refiz as contas de cada alternativa. (a) 4$\pi$R$^2$ = 1024$\pi$ dá R = 16, diâmetro 32 > 20: não cabe. (b) Com g = 18$\surd$3, a altura é 27 e o raio da esfera inscrita é 27/3 = 9; o diâmetro 18 > 16 impede a passagem pela abertura. O argumento da resolução está certo: a abertura é cercada de parede por todos os lados, então um disco de raio 9 centrado nela sempre encosta na parede. (c) a$\surd$3/2 = 20$\surd$3 dá a = 40. A altura da pirâmide, a/$\surd$2 = 20$\surd$2, bate com o enunciado. O raio da esfera inscrita é a$\surd$6/6 $\approx$ 16,33, diâmetro $\approx$ 32,66 > 20: não cabe. (d) 1331 = 121 + 121 + h$^2$ dá h = 33. O prisma de 11$\times$11$\times$33 passa em pé pela abertura de 16$\times$16 e ocupa exatamente a altura da caixa. A resposta é D, igual ao gabarito.%
\end{itemize}%
\paragraph{Pontos de dúvida}%
\begin{itemize}%
\item%
\textit{(statement)} Na alternativa (d), a altura do prisma (33 cm) é igual à altura da caixa indicada na figura. Isso só funciona se as paredes tiverem espessura desprezível ('caixa idealizada'), ou seja, se as dimensões internas forem iguais às externas. Nesse caso, a face superior do prisma fica rente ao plano da abertura, e entende{-}se que 'totalmente dentro' admite essa igualdade. A resolução declara essa hipótese, e ela é a leitura pretendida pela banca.%
\item%
\textit{(statement)} Na alternativa (c), 'altura 20$\surd$2' só é coerente com o apótema 20$\surd$3 se for a altura de cada pirâmide que forma o octaedro (a = 40), e não a distância entre vértices opostos. A resolução trata as duas leituras e chega a 'não cabe/não passa' em ambas. Na leitura alternativa, porém, a margem é pequena: 16,33 contra 16.%
\end{itemize}%
\section*{Questão 21}%
\begin{longtable}{@{}p{0.06\textwidth} p{0.33\textwidth} p{0.33\textwidth} p{0.1\textwidth} p{0.08\textwidth}@{}}\toprule Item & Resolução & Gabarito & Veredito & Certo \\ \midrule
21 & (B) 2000 & B & igual & ambos \\
\bottomrule\end{longtable}%
\begin{itemize}%
\item%
\textbf{21)} Refiz a conta. 12x$^2$ - 63y$^2$ = 5040 equivale a x$^2$/420 - y$^2$/80 = 1, que é uma hipérbole. Daí c$^2$ = 420 + 80 = 500 e F1O = c = 10$\surd$5. D e E são simétricos em relação a O e a reta DE é perpendicular ao eixo focal, então D e E ficam no eixo y. Com DE = 4c, o quadrilátero F1DF2E é um losango de diagonais 2c e 4c. Área = (2c$\cdot$4c)/2 = 4c$^2$ = 2000. Conferi também pela fórmula do cadarço (shoelace) em python, que deu 2000. A transcrição bate com o original (página 2 do PDF).%
\end{itemize}%
\paragraph{Pontos de dúvida}%
\begin{itemize}%
\item%
\textit{(statement)} O enunciado não diz se D e E pertencem à curva $\lambda$. Não podem pertencer, porque a hipérbole não corta o eixo y. A resolução observa isso de forma explícita, e o resultado não depende do ponto.%
\end{itemize}%
\section*{Questão 22}%
\begin{longtable}{@{}p{0.06\textwidth} p{0.33\textwidth} p{0.33\textwidth} p{0.1\textwidth} p{0.08\textwidth}@{}}\toprule Item & Resolução & Gabarito & Veredito & Certo \\ \midrule
22 & (B) todas as afirmativas são verdadeiras & D (apenas duas afirmativas são verdadeiras) & diferente & indefinido \\
\bottomrule\end{longtable}%
\begin{itemize}%
\item%
\textbf{22)} Refiz as contas e conferi com python3. P(x) = 3(x$^2$-x-2)(x$^2$+25), com raízes 2, -1, 5i e -5i (P se anula nas quatro). (I) Na convenção {[}0,2$\pi${[}, os argumentos são 0, $\pi$/2, $\pi$, 3$\pi$/2. Na convenção {]}-$\pi$,$\pi${]}, são -$\pi$/2, 0, $\pi$/2, $\pi$. Nos dois casos formam uma PA de razão $\pi$/2, então (I) é verdadeira. (III) $\sigma$$_{2}$ = c/a = 69/3 = 23, que é primo; a soma direta dos produtos das raízes também dá 23. (III) é verdadeira. A divergência está toda em (II). As áreas possíveis com vértices nos afixos são: quadrilátero convexo 15; triângulos \{2, -1, $\pm$5i\} 7,5; \{2, 5i, -5i\} 10; \{-1, 5i, -5i\} 5. O original da p. 3 confere com a transcrição: 'há um polígono cujo número que representa sua área {[}...{]} é divisor de 5'. Lido ao pé da letra, como existência, o triângulo de vértices -1, 5i e -5i tem área 5, que divide 5. Aí (II) é verdadeira e a resposta é B, como na resolução. O gabarito D só fecha se (II) for falsa, porque (I) e (III) são verdadeiras sem depender de convenção. Para isso, o polígono tem de ser o quadrilátero com os quatro afixos como vértices: área 15, que é múltiplo de 5, não divisor. Essa é claramente a leitura que a banca pretendia, com a pegadinha divisor $\times$ múltiplo. As contas da resolução estão certas e ela até calcula a área 15, mas escolhe a leitura existencial sem avisar que é ambígua. A redação ('há um') favorece a resolução; a intenção evidente da banca e o gabarito oficial favorecem D. Não dá para decidir só pela matemática: é caso de recurso.%
\end{itemize}%
\paragraph{Pontos de dúvida}%
\begin{itemize}%
\item%
\textit{(statement)} A afirmativa (II) é ambígua. 'Considerando os afixos de x1, x2, x3 e x4, há um polígono...' pode querer dizer (a) o polígono cujos vértices são os quatro afixos (quadrilátero, área 15, não divide 5, logo falsa e gabarito D) ou (b) existe algum polígono com vértices entre os afixos (o triângulo -1, 5i, -5i tem área 5, logo verdadeira e resposta B). O 'há um' favorece (b). A pegadinha divisor $\times$ múltiplo (15 é múltiplo de 5) indica que a banca pensou em (a). Pode ser motivo de recurso, anulação ou troca de gabarito; vale ver se a banca publicou o gabarito definitivo depois dos recursos.%
\item%
\textit{(key)} O gabarito D só se sustenta se (II) for falsa. (I) é verdadeira tanto em {[}0,2$\pi${[} quanto em {]}-$\pi$,$\pi${]}, e (III) dá $\sigma$$_{2}$ = 23, primo. Então o gabarito depende de ler (II) como 'o quadrilátero formado pelos quatro afixos', o que não está escrito no enunciado. Pelo jeito, a banca não percebeu que o triângulo -1, 5i, -5i tem área exatamente 5.%
\item%
\textit{(solution)} A resolução calcula a área do quadrilátero (15), mas escolhe a leitura existencial sem dizer que a afirmativa é ambígua e que a banca provavelmente quis o quadrilátero, o que levaria a D. Em prova de múltipla escolha, valeria comentar as duas leituras.%
\end{itemize}%
\section*{Questão 23}%
\begin{longtable}{@{}p{0.06\textwidth} p{0.33\textwidth} p{0.33\textwidth} p{0.1\textwidth} p{0.08\textwidth}@{}}\toprule Item & Resolução & Gabarito & Veredito & Certo \\ \midrule
23 & (B) 9 $\leq$ x < 13, com x = 11 mil dólares & B & igual & ambos \\
\bottomrule\end{longtable}%
\begin{itemize}%
\item%
\textbf{23)} Confere com o original (pág. 3 do PDF): a transcrição está fiel, inclusive a ordem trocada das colunas na 3ª tabela (Phenon, Eurocopter, Caravan). Refiz as contas em Python. Horas: Caravan 6$\cdot$1,5+14$\cdot$2+5$\cdot$3 = 52 h; Phenon 12$\cdot$1,5+4$\cdot$2+2$\cdot$3 = 32 h; Eurocopter 1$\cdot$1,5+8$\cdot$2+18$\cdot$3 = 71,5 h. Custos lidos pelo nome da aeronave: Caravan, na faixa 40--60 h, custa 4; Phenon, abaixo de 40 h, custa 3; Eurocopter, acima de 60 h, custa 4. Total x = 11, logo 9 $\leq$ x < 13, alternativa B. Se as colunas fossem lidas pela posição (a armadilha da questão), daria 4+1+8 = 13, alternativa C. A resolução aponta essa armadilha de forma explícita e correta.%
\end{itemize}%
\paragraph{Pontos de dúvida}%
\begin{itemize}%
\item%
\textit{(statement)} A faixa 'Entre 40 h e 60 h' não diz se os extremos 40 e 60 entram, e 'abaixo de'/'acima de' também não. Isso não afeta a resposta, porque nenhum total cai exatamente em 40 h ou 60 h (52; 32; 71,5). A resolução observa isso corretamente.%
\item%
\textit{(statement)} A ordem das colunas da 3ª tabela (Phenon, Eurocopter, Caravan) é diferente da 1ª (Caravan, Phenon, Eurocopter). É uma armadilha proposital: quem lê pela posição chega a x = 13, alternativa C. Não há ambiguidade, porque as colunas trazem o nome das aeronaves.%
\end{itemize}%
\section*{Questão 24}%
\begin{longtable}{@{}p{0.06\textwidth} p{0.33\textwidth} p{0.33\textwidth} p{0.1\textwidth} p{0.08\textwidth}@{}}\toprule Item & Resolução & Gabarito & Veredito & Certo \\ \midrule
24 & (A) 2 + P\_3$\cdot$C\_\{26,3\}$\cdot$A\_\{10,2\}$\cdot$C\_\{21,1\} = 29 484 002 & A & igual & ambos \\
\bottomrule\end{longtable}%
\begin{itemize}%
\item%
\textbf{24)} A transcrição confere com a página 10 do PDF (página 4 do arquivo). O teclado tem 26 letras, 21 símbolos e 10 algarismos. Refiz as contas em python. C\_\{26,3\}=2600 e A\_\{10,2\}=90. A alternativa a) vale 6$\cdot$2600$\cdot$90$\cdot$21+2 = 29 484 002, e b) vale 4 914 002. As alternativas c) e d) valem ambas 58 968 000, porque P\_3$\cdot$C\_\{26,3\}=A\_\{26,3\} e A\_\{21,1\}=C\_\{21,1\}. O raciocínio está correto. Os três modos de entrar se excluem, então somam{-}se 1+1+N\_c, o que elimina c) e d). Escolher as letras é uma combinação, o par de algarismos distintos é um arranjo e há 21 símbolos. O fator P\_3 vem da ordem dos três blocos, e não há dupla contagem porque os tipos de tecla são disjuntos. Essa é a leitura que bate com a) e com o gabarito oficial.%
\end{itemize}%
\paragraph{Pontos de dúvida}%
\begin{itemize}%
\item%
\textit{(statement)} O enunciado não diz se os três blocos (letras, algarismos, símbolo) têm posição fixa ou livre no código. O fator P\_3 da alternativa a) só se justifica se a ordem dos blocos for livre. Com formato fixo, a resposta seria b). O gabarito oficial (A) confirma a leitura de ordem livre adotada na resolução, mas um aluno poderia contestar a questão com base nessa ambiguidade.%
\item%
\textit{(statement)} A redação "um bloco de pares ordenados de algarismos distintos" está no plural, o que é ambíguo. A resolução lê como um único par, apoiada no fato de todas as alternativas trazerem A\_\{10,2\}. É razoável, mas é uma hipótese de interpretação.%
\item%
\textit{(solution)} A resolução supõe implicitamente que cada bloco é contíguo, isto é, que os tipos de tecla não se intercalam no código. Também supõe que o sistema aceita qualquer ordem de digitação das 3 letras, já que elas contam como conjunto. As duas hipóteses são necessárias para chegar exatamente a P\_3$\cdot$C\_\{26,3\}, mas a resolução não as declara. Também conta 1 possibilidade para a face e 1 para a digital, que é a leitura padrão.%
\end{itemize}%
\section*{Questão 25}%
\begin{longtable}{@{}p{0.06\textwidth} p{0.33\textwidth} p{0.33\textwidth} p{0.1\textwidth} p{0.08\textwidth}@{}}\toprule Item & Resolução & Gabarito & Veredito & Certo \\ \midrule
25 & (A) 8$\pi$(2 + $\surd$3/3) $\approx$ 64,8 u.v. & A & igual & ambos \\
\bottomrule\end{longtable}%
\begin{itemize}%
\item%
\textbf{25)} Conferi o enunciado no PDF original (pág. 4) e a transcrição está fiel. Refiz as contas. P(x) = -(x$^2$ - 2x + 4)(x$^2$ + 2x + 10), com raízes 1 $\pm$ i$\surd$3 e -1 $\pm$ 3i (testadas numericamente: |P(z)| = 0). Os afixos formam um trapézio isósceles com bases verticais 2$\surd$3 (em x = 1) e 6 (em x = -1) e altura 2. A rotação é em torno de x = 1, e o sólido é um cilindro (r = 2, h = 6) menos dois cones (r = 2, h = 3 - $\surd$3): V = 24$\pi$ - (8$\pi$/3)(3 - $\surd$3) = 16$\pi$ + 8$\pi$$\surd$3/3 = 8$\pi$(2 + $\surd$3/3) $\approx$ 64,776. Cheguei ao mesmo valor por três caminhos: cilindro menos cones, Pappus--Guldin com d = h(a+2b)/(3(a+b)) e integração numérica por arruelas. As outras alternativas dão cerca de 31,87 (b), 29,02 (c) e 13,73 (d). A resolução e o gabarito estão corretos e coincidem.%
\end{itemize}%
\section*{Questão 26}%
\begin{longtable}{@{}p{0.06\textwidth} p{0.33\textwidth} p{0.33\textwidth} p{0.1\textwidth} p{0.08\textwidth}@{}}\toprule Item & Resolução & Gabarito & Veredito & Certo \\ \midrule
26 & (D): curva logística decrescente, da assíntota 70 mil até perto de 0, com inflexão em (10, 35) mil, ou seja, P(t) = 70/(1+e\^{}\{k(t-10)\}), k > 0 & D & igual & ambos \\
\bottomrule\end{longtable}%
\begin{itemize}%
\item%
\textbf{26)} Conferi com a página 5 do PDF: a transcrição está fiel, inclusive o trecho "{]}-L, L{[}" e os quatro gráficos. A resolução está certa. A lei dada sempre dá valores em {]}0, L{[}, porque o denominador 1 + e\^{}\{\ldots{}\} é maior que 1. A curva vai de 0 a L, tem inflexão em (t0, L/2) e é simétrica em relação a esse ponto. Também refiz f'' = Lk$^2$u(u-1)/(1+u)$^3$ e está certo. Com isso, "amplitude de 70000" dá L = 70 mil, "a partir de 10 anos" dá t0 = 10 com P(10) = 35 mil, e "decrescimento" pede a curva decrescente. Isso corresponde ao gráfico (d), que mostra a assíntota 70 à esquerda, a curva tendendo a 0 à direita e o tracejado em t = 10 à meia altura. As alternativas (a) e (b) mudam de sinal e passam por zero em t0, o que nenhuma logística desse tipo faz, e população não pode ser negativa. A (c) é crescente. A resolução e o gabarito concordam.%
\end{itemize}%
\paragraph{Pontos de dúvida}%
\begin{itemize}%
\item%
\textit{(statement)} O enunciado afirma que "{]}-L, L{[} é o intervalo de mínimo e máximo" da curva. Isso contradiz a própria lei, cuja imagem é {]}0, L{[}, e o exemplo, em que L = 6 e a curva fica entre 0 e 6. Lendo ao pé da letra, amplitude 70 mil daria 2L = 70, ou seja, L = 35, e uma curva decrescente entre 35 e -35, que é exatamente o gráfico (a). A alternativa (a) parece ter sido feita para quem segue essa frase. Essa incoerência interna pode motivar recurso ou anulação. Mesmo assim, a fórmula, o exemplo e o contexto (população não pode ser negativa) levam a (d), como dizem a resolução e o gabarito.%
\item%
\textit{(statement)} O enunciado impõe k > 0, e com k > 0 a família f(t) = L/(1+e\^{}\{-k(t-t0)\}) é sempre crescente. Para modelar o decrescimento é preciso sair da família dada: trocar o sinal do expoente, como a resolução faz com g(t) = L - f(t), ou admitir k < 0. A resolução explica isso, mas o enunciado não.%
\item%
\textit{(solution)} A resolução lê "tendência a estabilização a partir de 10 anos" como ponto de inflexão em t0 = 10. A leitura é razoável: depois da inflexão, |f'| só diminui. Ela também coincide com o tracejado em t = 10 dos gráficos (c) e (d). Ainda assim, é uma interpretação do texto. Pela figura, a inflexão de (d) fica à meia altura, mais ou menos em 35, mas o valor 35 não está escrito no eixo.%
\end{itemize}%
\section*{Questão 27}%
\begin{longtable}{@{}p{0.06\textwidth} p{0.33\textwidth} p{0.33\textwidth} p{0.1\textwidth} p{0.08\textwidth}@{}}\toprule Item & Resolução & Gabarito & Veredito & Certo \\ \midrule
27 & B: Q3 < Q2 < Q1, com Q1 = 1, Q2 = 3(3$\surd$3-$\pi$)/(9$\surd$3-$\pi$) $\approx$ 0,495 e Q3 = (6$\surd$3-3$\pi$)/(6$\surd$3-2$\pi$) $\approx$ 0,235 & B & igual & ambos \\
\bottomrule\end{longtable}%
\begin{itemize}%
\item%
\textbf{27)} Refiz a questão olhando a figura original ampliada a 300 dpi. Q1: as duas figuras são retângulos 8r$\times$4r com área branca 8$\pi$r$^2$ (8$\cdot$$\pi$r$^2$ em cima e 2$\cdot$$\pi$(2r)$^2$ embaixo), então Q1 = 1. Q2: na figura de baixo o círculo inscrito está mesmo cinza e só os três círculos pequenos são brancos. Com VI = 2R e IJ = 2R-2$\rho$ = R+$\rho$, sai $\rho$ = R/3. Daí N2 = ?$^2$(3$\surd$3-$\pi$)/12, D2 = ?$^2$(9$\surd$3-$\pi$)/36 e Q2 $\approx$ 0,4952. Q3: cada círculo tem o centro no apótema, à distância a-$\rho$' do centro do hexágono. Centros vizinhos formam com o centro um triângulo equilátero, então a-$\rho$' = 2$\rho$' e $\rho$' = a/3 = m$\surd$3/6. Conferi numericamente que a distância entre centros vizinhos é igual a 2$\rho$'. Também confere com a figura: lado $\approx$ 195 px e diâmetro dos círculos $\approx$ 112 px. Q3 = (6$\surd$3-3$\pi$)/(6$\surd$3-2$\pi$) $\approx$ 0,2355. Portanto Q3 < Q2 < Q1, que é a alternativa B, igual ao gabarito. Conferi todos os valores intermediários com python3.%
\end{itemize}%
\paragraph{Pontos de dúvida}%
\begin{itemize}%
\item%
\textit{(statement)} A resposta depende de como se leem as cores na figura de baixo de Q2: lá o círculo inscrito é cinza. Ampliando a figura original, confirma{-}se que ele é cinza. Se fosse lido como branco, Q2 passaria de 1 e a ordem seria Q3 < Q1 < Q2, que não está entre as alternativas. As próprias alternativas, portanto, confirmam essa leitura.%
\item%
\textit{(statement)} Em Q3, o enunciado não diz em texto que cada círculo é tangente a um lado e aos dois vizinhos. Isso só se tira das marcas de ângulo reto na figura. A resolução declara essa hipótese, confere que os círculos ficam dentro do hexágono (distância 2$\rho$' > $\rho$' até os lados adjacentes), e as proporções da figura batem com $\rho$' = m$\surd$3/6.%
\end{itemize}%
\section*{Questão 28}%
\begin{longtable}{@{}p{0.06\textwidth} p{0.33\textwidth} p{0.33\textwidth} p{0.1\textwidth} p{0.08\textwidth}@{}}\toprule Item & Resolução & Gabarito & Veredito & Certo \\ \midrule
28 & (A): 33x + 242y - 1300 $\leq$ 0 e 5x + 7y - 91 $\geq$ 0 & A & igual & ambos \\
\bottomrule\end{longtable}%
\begin{itemize}%
\item%
\textbf{28)} Conferi o enunciado com a página 7 do PDF e a transcrição está fiel. Tiros: 165x + 231y $\geq$ 3003; dividindo por 33, fica 5x + 7y $\geq$ 91. Custo: 660x + 4840y $\leq$ 26000; dividindo por 20, fica 33x + 242y $\leq$ 1300. Os fatores conferem: 165=33$\cdot$5, 231=33$\cdot$7, 3003=33$\cdot$91, 660=20$\cdot$33, 4840=20$\cdot$242, 26000=20$\cdot$1300. Isso dá exatamente o sistema da alternativa (a). Também conferi as contas auxiliares da resolução: o par (19,0) dá 3135 tiros e custa R\$ 12540; 33$\cdot$19 - 1300 = -673; e 6,6$\cdot$91 = 600,6. Uma enumeração em Python achou 50 pares inteiros não negativos que atendem às condições. Todos satisfazem (a), nenhum satisfaz (c) ou (d), e todos também satisfazem (b), como a própria resolução aponta.%
\end{itemize}%
\paragraph{Pontos de dúvida}%
\begin{itemize}%
\item%
\textit{(statement)} Ambiguidade lógica do enunciado: ele diz que o par 'necessariamente deverá pertencer ao conjunto solução do sistema'. Com x, y $\geq$ 0, todo par que atende às condições também satisfaz a alternativa (b): 33x + 242y $\geq$ 6,6(5x + 7y) $\geq$ 600,6 $\geq$ 91 e 5x + 7y $\leq$ 33x + 242y $\leq$ 1300. Confirmei isso enumerando os 50 pares que atendem às condições. Lendo ao pé da letra, então, (a) e (b) são condições necessárias, e isso poderia sustentar um recurso pedindo a anulação. A intenção clara da banca, confirmada pelo gabarito oficial, é o sistema que modela o problema, que é o da (a). A resolução trata esse ponto de forma correta e transparente.%
\item%
\textit{(solution)} A resolução trabalha com x, y $\in$ $\mathbb{N}$ (x, y $\geq$ 0) sem que o enunciado diga isso. É uma hipótese natural, porque x e y são números de caixas, mas ela é usada no argumento sobre (b) e (d). Não afeta a escolha da alternativa (a).%
\end{itemize}%
\section*{Questão 29}%
\begin{longtable}{@{}p{0.06\textwidth} p{0.33\textwidth} p{0.33\textwidth} p{0.1\textwidth} p{0.08\textwidth}@{}}\toprule Item & Resolução & Gabarito & Veredito & Certo \\ \midrule
29 & (C): f(m) - m = $\surd$2 $\approx$ 1,414, que está em {[}1, 3/2{[} & C & igual & ambos \\
\bottomrule\end{longtable}%
\begin{itemize}%
\item%
\textbf{29)} Conferi com a figura da página 7: pontos (-9/4, 3), (-3/4, 1), (9/4, 5) e (15/4, 3). Daí saem m = 3, |t| = 2 e período 6, logo |q| = $\pi$/3. Escrevendo f(x) = 3 + 2cos($\pi$/3$\cdot$(x - 9/4)), a função passa pelos quatro pontos marcados, o que conferi numericamente. Então f(3) = 3 + 2cos($\pi$/4) = 3 + $\surd$2, e f(m) - m = $\surd$2 $\approx$ 1,414 $\in$ {[}1, 3/2{[}. A resolução está certa e chega à mesma alternativa do gabarito (C). O valor de m é a linha média, (1+5)/2, e não depende do sinal de t nem de q. Por isso o resultado não muda com a forma escolhida para escrever a função.%
\end{itemize}%
\paragraph{Pontos de dúvida}%
\begin{itemize}%
\item%
\textit{(statement)} Os valores 1, 3 e 5 e as abscissas -9/4, -3/4, 9/4 e 15/4 foram lidos dos rótulos da figura. Não há escala implícita nem leitura aproximada, e a figura é coerente: do mínimo ao máximo vai meio período (3), e dos extremos aos pontos na linha média vai um quarto de período (3/2).%
\end{itemize}%
\section*{Questão 30}%
\begin{longtable}{@{}p{0.06\textwidth} p{0.33\textwidth} p{0.33\textwidth} p{0.1\textwidth} p{0.08\textwidth}@{}}\toprule Item & Resolução & Gabarito & Veredito & Certo \\ \midrule
Questão 30 & C) 11 (g?(4), h?(6), j?(1)) & Questão anulada (qualquer alternativa aceita) & igual & ambos \\
\bottomrule\end{longtable}%
\begin{itemize}%
\item%
\textbf{Questão 30)} A banca anulou a questão e aceita qualquer alternativa, então a resposta C está de acordo com o gabarito. Refiz as contas com um exemplo, f(x) = -(x+1)$^2$ + 3 (x\_v = -1, y\_v = 3, c = 2). Para h: h(0) = 0, as raízes são 0 e 2x\_v, e o vértice é (x\_v, y\_v - c) = (-1, 1). Isso bate com o gráfico (6). Para j: as raízes são x$_{1}$ + c e x$_{2}$ + c, e o vértice é (x\_v + c, y\_v). Isso bate com o gráfico (1). O problema está em g = -c$\cdot$f: g(0) = -c$^2$ = -4 e g(x\_v) = -c$\cdot$y\_v = -6. Já o gráfico (4) marca -c = -2 e -y\_v = -3, que são valores de -f(x), e não de g. Como c > 1, nenhum dos seis gráficos é exatamente o de g. Isso explica a anulação. A banca provavelmente queria g?(4), h?(6) e j?(1), com soma 11 (alternativa C). Os distratores saem de g?(2), (3) ou (5), ou de j?(2). A resolução percebeu a diferença de escala na observação e escolheu a resposta pretendida, que é a interpretação mais razoável. Mesmo assim, a anulação está correta.%
\end{itemize}%
\paragraph{Pontos de dúvida}%
\begin{itemize}%
\item%
\textit{(statement)} Pelo enunciado, g(x) = -c$\cdot$f(x) com c > 1. Então g(0) = -c$^2$ e o vértice de g é (x\_v, -c$\cdot$y\_v). O gráfico (4), o único com concavidade para cima e raízes x$_{1}$ e x$_{2}$, marca -c e -y\_v, que valem para -f(x). Nenhum dos gráficos representa g de forma exata, e esse é muito provavelmente o motivo da anulação. A resolução afirma que '(4) é o único gráfico compatível'. Isso força a interpretação, porque com esses rótulos o gráfico (4) não é o de g.%
\item%
\textit{(solution)} A frase final diz que as outras alternativas 'exigiriam associar g a (2), (3) ou (5)'. Isso só vale mantendo h e j fixos. A alternativa d) 12, por exemplo, também sai de g?(4), h?(6), j?(2). O argumento simplifica demais, mas a resposta não muda.%
\item%
\textit{(statement)} Os gráficos (1) e (6) dependem de hipóteses tiradas da figura. Em (1), aparece x\_v + c > 0 e x$_{1}$ + c < 0. Em (6), a parábola deveria passar exatamente pela origem, porque h(0) = 0, e pela figura isso fica só aproximado. No gráfico (2), os rótulos -x$_{1}$ e -x$_{2}$ estão trocados (-x$_{2}$ < 0 < -x$_{1}$). Isso não afeta a associação pretendida.%
\end{itemize}%
\section*{Questão 31}%
\begin{longtable}{@{}p{0.06\textwidth} p{0.33\textwidth} p{0.33\textwidth} p{0.1\textwidth} p{0.08\textwidth}@{}}\toprule Item & Resolução & Gabarito & Veredito & Certo \\ \midrule
31 & c) aeronave 5 ($\sigma$ = $\surd$6,5 $\approx$ 2,55 h; o valor mais próximo é 3 h) & C & igual & ambos \\
\bottomrule\end{longtable}%
\begin{itemize}%
\item%
\textbf{31)} Conferi o gráfico no PDF original (pág. 9) e os valores lidos estão certos: 7, 10, 8, 5, 3, 5, 8, 2. A média é 6 h e a soma dos quadrados dos desvios é 52. Pela fórmula populacional, $\sigma$ = $\surd$6,5 $\approx$ 2,5495 h. A distância até 3 h (aeronave 5) é 0,450 e até 2 h (aeronave 8) é 0,550. Pela fórmula amostral, s $\approx$ 2,7255 h, e 3 h continua sendo o valor mais próximo (distância 0,274). Conferi as contas em Python. A resposta é c) nas duas fórmulas, o que bate com o gabarito.%
\end{itemize}%
\paragraph{Pontos de dúvida}%
\begin{itemize}%
\item%
\textit{(statement)} O enunciado não diz se o desvio padrão é populacional ou amostral. Isso não muda a resposta: dá $\approx$2,55 h numa e $\approx$2,73 h na outra, e as duas ficam mais perto de 3 h (aeronave 5). A resolução já confere esse ponto.%
\item%
\textit{(solution)} Os valores foram lidos no gráfico de barras. Conferi no original e as linhas tracejadas confirmam os números inteiros 7, 10, 8, 5, 3, 5, 8, 2. A resposta depende de $\sigma$$^2$ = 6,5 > 6,25. Se alguma barra fosse lida errado, a escolha entre 2 h e 3 h poderia mudar, mas não há ambiguidade na figura.%
\end{itemize}%
\section*{Questão 32}%
\begin{longtable}{@{}p{0.06\textwidth} p{0.33\textwidth} p{0.33\textwidth} p{0.1\textwidth} p{0.08\textwidth}@{}}\toprule Item & Resolução & Gabarito & Veredito & Certo \\ \midrule
32 & (C) 4 & C & igual & ambos \\
\bottomrule\end{longtable}%
\begin{itemize}%
\item%
\textbf{32)} A transcrição confere com o original (p. 9 do PDF, numerada 15 na prova). Com P(macho) = 2/5 e sorteios independentes, P(pelo menos uma fêmea) = 1 - (2/5)\^{}n > 0,95, o que equivale a (2/5)\^{}n < 1/20. Com os logaritmos dados, -0,4n < -1,3, logo n > 3,25 e o mínimo é n = 4. A conferência direta bate: (2/5)\^{}3 = 0,064 > 0,05 (93,6\%, não serve) e (2/5)\^{}4 = 0,0256 < 0,05 (97,44\%, serve). Com logaritmos exatos (log 0,4 $\approx$ -0,398 e log 0,05 $\approx$ -1,301), dá n > 3,27, e a conclusão continua a mesma. A resolução e o gabarito concordam e estão certos.%
\end{itemize}%
\paragraph{Pontos de dúvida}%
\begin{itemize}%
\item%
\textit{(statement)} O enunciado fala da proporção de nascimentos (2 machos para 3 fêmeas), não da composição atual do rebanho, e a amostra é tirada sem reposição de um rebanho finito. A resolução supõe, e declara isso, que o rebanho é grande e que os sexos são independentes com P(macho) = 2/5. É a leitura que a banca quis, já que dá log(2/5). Num rebanho pequeno, o modelo hipergeométrico poderia fazer n = 3 servir, mas não há dados para isso.%
\end{itemize}%
\end{document}